\chapter{从直线到曲面：局部近似的预告}

\section{一场雨中的自行车赛}

环青海湖自行车赛进行到下坡路段，天上下起了小雨。电视镜头从直升机上俯拍：盘山公路像一条灰色的丝带，在绿山坡上绕出一个又一个弯。解说员反复念叨一句话：``前面这个弯很急，车手一定要提前减速。''

镜头里，车手们的表现确实不一样。经过一些弯道时，他们几乎不用捏刹车，身体轻轻一斜就过去了；经过另一些弯道时，他们早早把速度降下来，车身倾斜得很厉害，轮胎在湿漉漉的路面上划出一道弧线。显然，``急''和``缓''是有客观差别的，不只是解说员的修辞。

现在请你当一回赛事工程师。组委会给你一张赛道的平面图——就是地图上那种弯弯曲曲的线，没有等高线，没有比例尺之外的任何信息。他们要你回答两个听起来很简单的问题：

第一个问题：图上这些弯，哪个最急？你能给``弯得有多厉害''配一把尺子吗？

第二个问题：为什么这张地图放大很多倍之后，再弯的路，每一小段看起来都几乎是直的？

先别往下看，凭直觉猜一猜。对于第一个问题，大多数人会说``转弯转得快就是急''——这话没错，但它把问题换了个说法而已：什么叫``转得快''？用什么单位？对于第二个问题，很多人会愣一下：这好像是天经地义的，可一旦追问``为什么''，就发现它并不像看起来那么理所当然。这两个问题，正是本章的起点。

\section{第一次尝试：几根不够格的尺子}

要量``弯''，最朴素的想法是借用直线。曲线我们不熟悉，直线我们太熟悉了，那就看看这条弯道``离直线有多远''。

第一根尺子：连接弯道起点和终点，量这条线段的长度。弯道越急，两端点应该被``甩''得越近，线段越短，对吗？不对。找一条长长的、平缓的大弧线，再找一个短短的、突然掉头的小急弯，你完全可以让它们的端点距离一模一样。端点距离反映的是``你取的这段有多长''，而不是``这段弯得有多急''。

第二根尺子：还是连接起点和终点，这次量曲线离这条弦最远有多远，也就是所谓的``弓高''。弓越高越弯？还是不行。同一个弯，你取一大段来看，弓可以很高；只取中间一小段，弓又矮下去。尺子随着你取的区间忽大忽小，这样的量法没有资格叫``程度''——程度应该是弯道自己的性质，不该由旁观者随便改变。

两次失败透露了一个重要线索：``弯曲''必须是一个在一点附近说了算的性质。你不能站在远处、圈一大段来评判一个弯急不急；你必须凑近了看，在任意小的一段上，它的弯曲程度都应该有确定的答案。数学家把这种``在一点附近''的视角叫做\textbf{局部}视角。这一章的核心思想只有一句话：弯曲的东西，在局部是直的；而研究弯曲的办法，就是先研究它在局部像什么样的直线、又怎样一点一点偏离这条直线。

可是``凑近了看''到底是什么意思？凑多近才算近？我们需要一个更清楚的图景。

\section{把曲线放到显微镜下}

先看一个你早就知道、却未必细想过的例子：地球是圆的，可你脚下的操场看起来是平的。这不是因为操场特殊，而是因为你太小、地球太大。我们来算一笔账。假设你沿一个半径为 $R$ 的圆周走，每走一段弧长 $s$，你的前进方向改变的角（用弧度制）是 $s/R$。地球半径约 $6\,400\,000$ 米，你沿地面走 $1$ 米，方向只改变
\[
\frac{1}{6\,400\,000}\ \text{弧度} \approx 0.000\,008\,9^\circ.
\]
这个角度小到任何量角器都量不出来，于是你的全部生活经验都在告诉你：地面是平的。古人认为大地是平的，不是因为他们笨，而是因为``局部像平面''这件事实在太逼真了。反过来，当人们终于掌握了足够精确的测量手段——比较两地正午影子的长度、观察远航船只先消失船身再消失桅杆——大地的弯曲才从``局部的假象''里浮现出来。局部的平直与整体的弯曲并不矛盾，它们是同一件事物在不同尺度上的两副面孔。

把同一个想法倒过来用：取一条弯弯曲曲的曲线，盯住它的一小段，放大；再放大；再放大。只要这条曲线足够``顺''——不折断、不起尖角——你会发现，放大到足够程度之后，它看起来就是一条直线。圆如此，抛物线如此，盘山公路的地图也如此。这不是巧合，而是``光滑''二字的含义：光滑曲线在显微镜下是直的。

\begin{fanli}
``任何曲线，只要放得足够大，看起来都像直线。''这句话是错的。请画一个字母 V，也就是函数 $y=|x|$ 的图像，盯住它的尖角放大。放大一倍、十倍、一百倍——尖角还是那个尖角，张开的度数一丝一毫都没有变。放大改变的是尺寸，不是角度。所以``局部像直线''是光滑曲线才有的特权；有尖角、有断裂的曲线不享受这个待遇。数学家把``局部像直线''这条性质起了一个正式的名字，叫``可微''，它将是第三部的主角之一。现在你只要记住：本章谈的都是没有尖角的光滑曲线。
\end{fanli}

\begin{shiyan}
不用任何仪器，你就能亲眼看到``放大变直''。在纸上随手画一条弯弯曲曲的线，用铅笔在其中一段上涂一个指甲盖大小的记号。把纸拿到复印机上，以记号为中
心放大复印一次；把复印件再放大复印一次。两三次之后，记号附近的那段线在复印件上几乎就是一条直线段。如果家里有放大镜，直接用它看曲线上的一小段，效果也一样。然后换一张画着字母 V 的纸重复这个过程，对比尖角处的表现——反例就在你手里。
\end{shiyan}

\section{切线：曲线在一点处的替身}

既然光滑曲线的一小段像直线，下一个动作顺理成章：把这条直线找出来，让它当曲线在这一点附近的``替身''。直线我们太熟悉了——一个方向、一个斜率，一个数就能说完。如果弯曲的问题能翻译成直线的问题，难度立刻降一个档次。

可是，``最贴切的那条直线''是哪一条？这里有一个陷阱。你可能会说：切线嘛，就是只碰曲线一个点的直线。这个印象来自圆，但它搬到一般曲线身上立刻出问题。抛物线 $y=x^2$ 的顶点处，竖直的那条直线也只碰抛物线一个点，可它显然不是我们想要的``替身''——它跟抛物线的走势完全对不上。反过来，一条在局部最贴切的直线，放远了看完全可能再碰到曲线别的地方。``只碰一个点''是圆的切线的性质，不是切线的定义。

正确的想法来自上一节的失败与这一节的放大：盯住曲线上一点 $P$，在旁边再取一点 $Q$，连成一条直线 $PQ$。这样的直线叫\textbf{割线}，它切进曲线，告诉我们从 $P$ 到 $Q$ 这段时间里曲线的``平均走向''。现在让 $Q$ 沿着曲线慢慢滑向 $P$：割线随之转动，转过的角度越来越小，位置越来越稳定，最后稳定下来的那个位置，就是曲线在 $P$ 点的\textbf{切线}。

\begin{center}
\begin{tikzpicture}[scale=1.1]
  \draw[->] (-0.5,0) -- (4.5,0) node[right] {$x$};
  \draw[->] (0,-0.5) -- (0,4.2) node[above] {$y$};
  \draw[thick] plot[smooth] coordinates {(-0.3,0.05) (0.5,0.15) (1,0.5) (1.8,1.6) (2.6,3.2) (3.2,4.0)};
  \fill (1,0.5) circle (1.5pt) node[below left] {$P$};
  \fill (2.6,3.2) circle (1.5pt) node[above right] {$Q_1$};
  \fill (1.8,1.6) circle (1.5pt) node[above right] {$Q_2$};
  \draw (0.4,-0.1) -- (3.4,3.9) node[right] {割线 $PQ_1$};
  \draw[dashed] (0.2,-0.4) -- (3.0,2.8) node[right] {割线 $PQ_2$};
  \draw[thick] (-0.2,-0.7) -- (3.4,3.05) node[right] {切线};
\end{tikzpicture}
\end{center}

\begin{dingyi}[切线的直观定义]
设 $P$ 是光滑曲线上一点。在曲线上取另一点 $Q$，当 $Q$ 沿曲线无限接近 $P$ 时，割线 $PQ$ 趋向的那条确定直线，称为曲线在 $P$ 点的切线。
\end{dingyi}

注意这个定义里藏着一个狡猾的东西：$Q$ 永远不能真的等于 $P$——两点才能定一条直线，$Q$ 一旦撞上 $P$，割线就不存在了。切线是由一个``永远走不到头的过程''决定的。现在我们先用直觉接受这个过程；它引起的麻烦，以及数学家怎样把它驯服成严格的``极限''，是下一章的故事。

\begin{shuyu}{切线}
直观解释：曲线在一点处的``最佳直线替身''，显微镜下曲线呈现的样子。正式定义：割线 $PQ$ 在 $Q$ 趋向 $P$ 时的极限位置。一个例子：圆上任意一点的切线，垂直于该点处的半径。
\end{shuyu}

最后一个例子值得细看，因为它能检验新想法和老知识对不对得上。初等几何告诉我们：圆的切线垂直于半径。用割线逼近的眼光能重新看出这件事吗？能。设圆心为 $O$，在圆上取点 $P$，再取两个关于直线 $OP$ 对称的点 $Q$ 和 $Q'$。由对称性，弦 $QQ'$ 垂直于 $OP$。现在让 $Q$ 和 $Q'$ 同时沿圆周滑向 $P$：每一步它们仍然对称，割线仍然垂直于 $OP$；于是极限位置——切线——也垂直于 $OP$。新方法没有推翻老结论，而是把老结论变成了新图景里的一个特例。一个数学想法靠不靠谱，这是重要的试金石。

切线这把``替身''有什么用？用处大了。车轮沿弯道滚动时，如果突然打滑，车会沿切线方向飞出去——链球运动员松手后，链球正是沿切线飞出的；砂轮磨刀具时溅出的火花，也是沿切线射出的。曲线在每个瞬间的``意图''，都写在它的切线上。

\section{弯得有多厉害：曲率}

有了切线，``方向''这个词就精确了：曲线在一点的方向，就是切线的方向。现在可以正式打造量``弯''的尺子。

想象你是一只沿曲线爬行的蚂蚁，身上装着指南针，随时记录自己前进方向的角度。走同样长的一段路：在平缓的地方，指南针几乎不动；在急弯处，指南针呼呼地转。于是弯曲的程度可以定义为：单位弧长上，前进方向改变了多少。

\begin{dingyi}[曲率的初步定义]
光滑曲线上一点处的曲率，是切线方向角沿弧长的变化率：在这一点附近，每走过一单位弧长，切线方向角改变的弧度数。曲率越大，弯得越急。
\end{dingyi}

这个定义有点抽象，好在有一类曲线可以把它算到底——圆。

\begin{dingli}[圆的曲率]
半径为 $R$ 的圆，其曲率处处相同，等于 $\dfrac{1}{R}$。
\end{dingli}

\begin{proof}
在圆周上取一点，设从它出发沿圆周走过弧长 $s$，到达另一点。这段弧对应的圆心角记为 $\theta$（用弧度制）。由弧长公式，
\[
s = R\theta.
\]
再看方向改变了多少。半径始终垂直于切线；半径从起点转到终点，转过了圆心角 $\theta$，而``垂直''这个关系在转动中保持不变，所以切线方向也恰好转过 $\theta$。于是方向改变量与弧长之比为
\[
\frac{\theta}{s} = \frac{\theta}{R\theta} = \frac{1}{R}.
\]
这个比值里不含有起点的任何信息，所以它在圆周上处处相同。这就是我们要证的。
\end{proof}

\begin{li}
一个半径 $10$ 米的环形路口，沿它每走 $1$ 米，方向改变 $\dfrac{1}{10}$ 弧度，约 $5.7^\circ$；一个半径 $2$ 米的急弯，每走 $1$ 米方向要改变 $\dfrac{1}{2}$ 弧度，约 $28.6^\circ$。后者曲率是前者的五倍——弯越急，半径越小，曲率越大。直觉和公式对上了。
\end{li}

\begin{center}
\begin{tikzpicture}[scale=0.9]
  \draw[thick] (0,0) circle (1);
  \draw[->,thick] (1,0) -- (1,1.2);
  \node at (-1.6,1.1) {半径 $1$：曲率 $1$};
  \draw[thick] (5.5,0) circle (2.5);
  \draw[->,thick] (8,0) -- (8,1.2);
  \node at (5.5,3.1) {半径 $2.5$：曲率 $0.4$};
\end{tikzpicture}
\end{center}

\begin{zhu}
直线不转弯，方向角永远不变，曲率是 $0$。按公式 $\dfrac{1}{R}$ 理解，直线相当于``半径无穷大的圆''。这不是文字游戏，而是一个有用的视角：直线是弯曲程度为零的曲线，就像 $0$ 是``没有''的数一样自然。
\end{zhu}

\begin{shuyu}{曲率}
直观解释：弯得有多厉害——每走一单位路程，方向转过多少。正式定义：切线方向角沿弧长的变化率；圆的曲率等于半径的倒数。一个例子：半径 $4$ 厘米的圆，曲率是每厘米 $\dfrac{1}{4}$ 弧度。
\end{shuyu}

曲率不是书房里的摆设，工程师天天在和它打交道。铁路由直道接入弯道时，绝不能让曲率从 $0$ 突然跳成 $\dfrac{1}{R}$——列车会被这记``急刹车式的转向''狠狠甩一下。聪明的做法是在直道和圆弯之间夹一段过渡，让曲率从零开始连续地、缓缓地升到 $\dfrac{1}{R}$，乘客几乎感觉不到。高速公路的匝道、过山车的轨道、过山车回环入口处越来越紧的弧线，背后都是同一条设计原则：让``弯得有多厉害''这个量平滑地变化。你看，一个``单位弧长上方向改变多少''的抽象定义，最后落进了每个人的安全带里。

一般的曲线在不同地方曲率不同：盘山公路的直段曲率是 $0$，缓弯处小，急弯处大。数学家的处理办法漂亮而自然：在曲线上每一点，找一个与该点贴合得最好的圆——既不偏松也不偏紧——叫做这一点的\textbf{密切圆}，规定这点的曲率就是密切圆半径的倒数。直线处，密切圆半径无穷大；急弯处，密切圆又小又紧。这句话我们只当作直观的描述，它的严格定义需要导数，要等下一部才能补齐。但请注意我们做的事：用直线逼近曲线，再用圆逼近曲线``偏离直线的程度''——一步一步，把弯曲翻译成越来越精细的局部近似。这条``局部近似''的路，会一直通到微分几何的深处。

\section{换一个视角：球面上的``直线''}

到目前为止，我们都在平面上打转。现在换一个舞台，看看这些想法能走多远。

打开航班追踪软件，看一趟从北京飞往纽约的飞机。在常见的平面地图上，它的航迹不是一条直线，而是一条向北鼓起的大弧线，擦着北极圈过去。飞机为什么要绕远？答案是：它没有绕远。地图是把球面硬压平得到的，压平的过程扭曲了形状；在真实的球面上，这条``弯''的航线恰恰是几乎最直的路线。

平面上，``直''和``最短''是同一件事的两面：直线段是两点间最短的连接。把这个品格搬到球面上，就能找到球面世界里的``直线''：球面上两点间最短的路线。它们是什么？是用过球心的平面切出来的圆——\textbf{大圆}。赤道是大圆，所有经线圈是大圆；而赤道以外的纬线圈不是大圆，它们的平面不过球心，所以沿纬线飞并不划算。

\begin{shuyu}{大圆}
直观解释：球面上的``直线''，球面上最直也最短的路线。正式定义：球面与过球心的平面相截得到的圆，也称测地线。一个例子：赤道和每一条经线圈都是大圆；北京到纽约的短航线沿着大圆走。
\end{shuyu}

大圆登场之后，一件惊人的事发生了。在球面上，用大圆弧当``边''，可以围成三角形，叫做球面三角形。请看一个具体的：从北极点出发，沿一条经线向南走到赤道，再沿赤道向东走四分之一个赤道周长，再沿另一条经线向北回到北极。三条边都是大圆弧，是一个标准的球面三角形。算算它的内角和。

\begin{dingli}[有三个直角的三角形]
上述球面三角形的三个内角都是 $90^\circ$，内角和为 $270^\circ$。
\end{dingli}

\begin{proof}
分两步。第一步，看赤道上的两个角。在交点处，赤道的方向是正东，经线的方向是正北——东西方向与南北方向互相垂直，所以经线与赤道的夹角是 $90^\circ$，两个底角都是直角。第二步，看北极处的顶角。两条经线在北极的夹角等于它们的经度差；我们沿赤道走了四分之一个周长，对应的经度差恰好是 $360^\circ \div 4 = 90^\circ$。三个角都是 $90^\circ$，加起来 $270^\circ$。
\end{proof}

\begin{center}
\begin{tikzpicture}[scale=1.0]
  \draw (0,0) circle (2.6);
  \draw (0,0) ellipse (2.6 and 0.8);
  \draw[very thick] (0,2.6) .. controls (1.3,2.2) and (2.4,1.0) .. (2.6,0);
  \draw[very thick] (0,2.6) .. controls (-0.45,2.0) and (-0.45,1.3) .. (0,0.8);
  \draw[very thick] (2.6,0) arc[start angle=0, end angle=90, x radius=2.6, y radius=0.8];
  \fill (0,2.6) circle (1.5pt) node[above] {北极};
  \node at (1.35,0.35) {球面三角形};
\end{tikzpicture}
\end{center}

$270^\circ$！平面几何里，三角形内角和永远是 $180^\circ$，这是我们学过的最铁的定律之一。可在球面上，它失效了——不是差一点半点，而是整整多出 $90^\circ$。多出来的这部分叫\textbf{角盈}。先看看这个三角形到底有多大：两条经线把球面从北极到南极切成四瓣，赤道再把每瓣切成南北两半，我们的三角形正好占其中一块，也就是整个球面的八分之一。而且三角形越大，角盈越大：上面这个占了球面八分之一的大家伙角盈是 $90^\circ$；如果你沿赤道只走八分之一周长就回家，三角形变小，顶角变成 $45^\circ$，内角和变成 $225^\circ$，角盈也随之减半。十九世纪初，数学家把这份观察凝成了一条精确的定理：球面三角形的角盈与它的面积成正比（在半径为 $R$ 的球面上，角盈换算成弧度后恰好等于面积除以 $R^2$）。我们不给完整证明，但直观是清楚的：三角形越大，它``裹住''的弯曲越多；足够小的三角形几乎躺在平面上，内角和也就几乎回到 $180^\circ$。

请你停下来体会最后这句话的分量。一只生活在球面上的蚂蚁，不需要飞到太空去``看''地球是弯的——它只要在地面画一个三角形、量一量内角和，就能发现自己世界的弯曲。弯曲不只是一种``从外面看到的形状''，它写进了曲面内部的几何定律里，用尺子、量角器和光就够探测。数学家把这叫做弯曲的\textbf{内蕴}性质。我们脚下的大地、我们头顶的星空，都可以这样从内部探测。

\begin{lianjie}
本章的概念在书里书外都有远亲。向前看：``割线趋向切线''里的极限过程，是第 9 章的主角；``切线的斜率''正是第 10 章导数的几何面孔；``曲率''和``球面上的几何''将在第 17 章长成微分几何。向外看：全球卫星定位系统计算两点距离时用的是大圆而不是地图直线；爱因斯坦的广义相对论更进一步，认为引力不是力，而是时空本身的弯曲——行星绕太阳运行，走的正是弯曲时空里的``直线''。甚至连本章的蚂蚁都是真实的：测绘工作者正是靠在地面上测量三角形的内角和，来检验地球表面的形状。
\end{lianjie}

\begin{shiyan}
找一个橙子（或苹果、网球）、一支记号笔和一把软尺。在橙子中部画一个圆圈当``赤道''，再画两条经过``两极''的经线，让它们在极点处的夹角大约是直角——三条线围出一个球面三角形。把一小片纸撕成楔形当角度样板，或者直接用量角器贴着表面估一估三个角：你会发现每个角都接近 $90^\circ$，内角和明显大于 $180^\circ$。最后做一件更有破坏性的事：把橙子皮剥下来，试着把一块皮在平桌上摊平。你会发现它要么翘起来，要么裂开缝——球面无法既不拉伸又不撕裂地摊成平面。这正是所有世界地图都必须变形的原因，也是``弯曲''最实在的证据。
\end{shiyan}

\begin{toolbox}
本章往你的背包里放了五件新工具。
\begin{itemize}
  \item \textbf{局部近似（放大法）}：研究弯曲，先盯住一小段放大；光滑曲线在局部像直线，光滑曲面在局部像平面。
  \item \textbf{割线到切线}：用平均走向逼近瞬时走向，切线是割线的极限位置，是曲线在一点处的直线替身。
  \item \textbf{曲率}：单位弧长上方向角的改变量，量``弯得有多厉害''；圆的曲率 $= 1/R$，直线曲率为 $0$。
  \item \textbf{大圆（测地线）}：曲面上``最直''的路线，由``最短''来定义，球面上就是过球心平面截出的圆。
  \item \textbf{角盈}：球面三角形内角和超过 $180^\circ$ 的部分，与面积成正比；它让弯曲成为可以从曲面内部测量的性质。
\end{itemize}
\end{toolbox}

\begin{pandeng}
想继续往上爬的同学，可以记下这些关键词，日后在高中、大学或竞赛读物里与它们重逢：极限与导数（``无限接近''的严格化）；密切圆与曲率半径（曲线的二阶局部近似）；高斯曲率与高斯的``绝妙定理''——曲面的弯曲完全由曲面内部的测量决定；测地线与黎曼几何；以及广义相对论中``弯曲时空''的科普读物。每一个词背后，都站着本章已经见过的直观图景。
\end{pandeng}

\section*{思考题}

\textbf{观察题}

\begin{sikao}
  \item 找一段你熟悉的弯路（上学的路线、操场的弯道、游乐场的过山车图），标出你认为曲率最大和最小的地方。提示：想象自己开车经过，方向盘在哪里打得最急？哪里几乎不用动？
  \item 在地球仪上观察：赤道、经线圈、北纬 $60^\circ$ 纬线圈，哪些是大圆？沿北纬 $60^\circ$ 线从北京``正东''飞到同纬度的欧洲城市，是最短路线吗？提示：想象用一个平面去切地球，切面要经过球心；再看看这条纬线所在的平面经过球心吗？
\end{sikao}

\textbf{证明题}

\begin{sikao}
  \item 不用弧长公式 $s=R\theta$，改用``半圆''整体验证圆的曲率：沿半个圆走一圈，切线方向总共改变了多少？走过的路程是多少？两者相除，是否仍得 $\dfrac{1}{R}$？提示：半圆两端的切线方向恰好相反，圆弧长度是 $\pi R$。
  \item 证明本章那个``三个直角的球面三角形''的面积恰好是球面的八分之一，并验证：它的角盈（换算成弧度）等于面积除以 $R^2$。提示：整个球面的面积是 $4\pi R^2$；八分之一球面由四个这样的三角形拼成还是八个？角盈是 $270^\circ - 180^\circ$，注意把角度换成弧度。
\end{sikao}

\textbf{挑战题}

\begin{sikao}
  \item 把一张长方形的纸卷成圆筒（不拉伸、不撕破）。纸上原来的一条直线段，卷起来后变成了什么？圆筒面上的``三角形内角和''大于 $180^\circ$ 吗？提示：纸既没拉伸也没撕裂，说明纸面上的所有长度、角度都没有改变——一只生活在纸面上的蚂蚁，能用本章的内角和方法区分``平面''和``圆筒面''吗？这件事和球面的情况有什么本质不同？
  \item 一只蚂蚁带着一支小箭沿球面三角形的边爬行，每走一条边都让箭头尽量``指向前方、不打转''，转弯处随三角形的角改变方向。爬完一圈回到起点时，箭头的方向和出发时一样吗？差多少？提示：先在平面上试（答案应当是一样），再在球面直角三角形上试，把每次转弯的角度加起来，和角盈比一比。
\end{sikao}

\section*{通往下一章的门}

回头清点一下，本章我们其实只做了一件事：把``无限小''请进了数学。光滑曲线放大到无穷倍像直线；切线是割线在 $Q$ 无限接近 $P$ 时的归宿；曲率是``无限小一段''上方向的改变率；球面上足够小的三角形，内角和无限接近 $180^\circ$。每一个关键概念，都站在一个``无限接近却永不到达''的过程上。

可是这个过程，我们今天全是靠直觉接住的。``$Q$ 滑向 $P$''——滑到多近才算数？``无限接近''和``到达''到底差在哪？如果一套理论的地基建在这么模糊的动词上，我们凭什么相信它推出的结论？

两百多年前，数学家们就被同样的问题拷问着，并且吵了一个多世纪才给出答案。那个答案叫做\textbf{极限}——它能把``无限接近''翻译成不含糊的、可以检验的不等式语言。它是整个第三部的地基，也是导数、积分和一切``变化之学''的起点。下一章，我们就去直面这个躲在所有光滑与弯曲背后的幽灵：无限小，究竟是不是一个合法的数学对象？
